\documentclass[11pt]{article}

% ---------- Packages ----------
\usepackage[margin=1in]{geometry}
\usepackage{graphicx}
\usepackage{amsmath}
\usepackage{amssymb}
\usepackage{booktabs}
\usepackage{cite}
\usepackage{hyperref}
\usepackage{enumitem}

% ---------- Title, Authors ----------
\title{Failure Modes in Chained Tool-Calling Agents Operating Without Human Checkpoints}
\author{Rupesh Kumar\thanks{Department of Information Technology (AI \& ML), Indian Institute of Information Technology, Allahabad.}}
\date{}

\begin{document}

\maketitle

\begin{abstract}
Large language model (LLM)-based agents are increasingly being designed to perform
tasks through repeated interactions with external tools such as search engines,
databases, APIs, code interpreters, browsers, file systems, and enterprise services.
Unlike conventional language-model applications that primarily generate text,
tool-calling agents can transform model outputs into actions that modify external state.
When several tool calls are chained without human checkpoints, a small error in
planning, tool selection, parameter generation, interpretation of tool output, or state
tracking can propagate through subsequent steps and produce increasingly severe
failures. This paper examines the major failure modes associated with autonomous
chained tool calling, focusing on error propagation, incorrect tool selection, parameter
and schema errors, stale or misleading tool outputs, planning failures, state
inconsistency, prompt injection, privilege escalation, irreversible actions, and
inadequate recovery. Research on ReAct, Toolformer, API-Bank, Gorilla, AgentBench,
InjecAgent, AgentDojo, ToolEmu, and Agent Security Bench demonstrates that tool
use introduces capabilities as well as new reliability and security risks. The paper
argues that chained execution changes the risk profile of an agent because
downstream actions may treat an earlier model-generated assumption as an
established fact. Human checkpoints therefore represent more than simple approval
steps: they can function as control boundaries that interrupt error propagation and
prevent uncertain plans from becoming irreversible actions. Current research
increasingly emphasizes sandboxing, capability restriction, runtime monitoring,
structured tool interfaces, adversarial evaluation, and dynamic benchmarks. However,
significant gaps remain in measuring cumulative failure probability across long tool
chains and in determining when an agent should autonomously continue, recover,
request clarification, or stop. Future research should therefore combine reliability
engineering, security controls, provenance tracking, uncertainty estimation, and
adaptive human oversight into unified frameworks for safe agentic execution.
\end{abstract}

\noindent\textbf{Keywords:} large language models, AI agents, tool calling, function calling,
autonomous agents, failure propagation, prompt injection, human oversight,
agent security, reliability

\section{Introduction}

Large language models have evolved from systems that primarily generate
natural-language responses into systems capable of interacting with external
environments. Tool-augmented agents can retrieve information, execute calculations,
manipulate files, browse websites, call APIs, write and execute code, and perform
multi-step tasks. ReAct demonstrated an influential approach in which reasoning and
acting are interleaved, allowing an LLM to use observations from an external
environment to update its subsequent actions \cite{yao2023react}. Toolformer similarly
demonstrated that language models can learn when to invoke APIs, what arguments
to provide, and how to incorporate tool outputs into subsequent generation \cite{schick2023toolformer}.

The increasing autonomy of these systems creates an important distinction between
single tool use and chained tool use. In a single interaction, an incorrect tool call
may produce one incorrect result. In a chain, however, the output of tool $T_1$ can
become an input to $T_2$, whose output becomes an input to $T_3$, and so on. The
resulting process can be represented as:
\begin{equation}
S_0 \rightarrow A_1 \rightarrow O_1 \rightarrow A_2 \rightarrow O_2
\rightarrow \cdots \rightarrow A_n \rightarrow O_n
\label{eq:chain}
\end{equation}
where $S_0$ is the initial state, $A_i$ is an agent action or tool call, and $O_i$ is
the resulting observation. If an early action is inappropriate or its output is
misinterpreted, later actions may operate on a corrupted state or incorrect
assumption.

This problem becomes particularly significant when there are no human
checkpoints between tool calls. A human checkpoint is a deliberate point at which a
person reviews a proposed action, verifies an important intermediate result, approves
a high-impact operation, or resolves ambiguity before execution continues. Without
such checkpoints, the agent may autonomously transform uncertain reasoning into
real-world actions.

Existing benchmarks already demonstrate weaknesses in tool use, planning,
and agent security. API-Bank evaluates planning, API retrieval, and API calling and reports substantial remaining
challenges in tool utilization. AgentBench evaluates agents across multiple interactive
environments and identifies long-term reasoning, decision-making, and instruction-following
as major obstacles.

This paper examines the failure modes that arise when tool calls are chained together without
intermediate human validation. The central argument is that the principal risk is not merely that an
individual tool call can fail. Rather, errors can become inputs to subsequent decisions,
allowing local failures to propagate into compound failures.

\section{Background and Related Work}

\subsection{Tool-Augmented Language Models}

Tool use provides language models with capabilities that are difficult to obtain
through language generation alone. Toolformer showed that external APIs can
provide calculators, search systems, translation services, and other specialized
capabilities.

Gorilla further investigated API invocation and highlighted problems including
inaccurate arguments and hallucinated API usage. Its evaluation emphasized the
importance of accurate API selection and parameter generation, especially when the
number of available APIs is large or documentation changes over time \cite{patil2023gorilla}.

A later survey by Qu et al. organizes tool learning into four broad stages: task
planning, tool selection, tool calling, and response generation. This decomposition
is useful for understanding chained failures because an error at any stage can
affect every later stage \cite{qu2025survey}.

\subsection{Agent Evaluation}

API-Bank provides a structured benchmark for tool-augmented LLMs, including
73 runnable APIs and hundreds of annotated tool-use dialogues. Its analysis
shows that tool utilization involves multiple distinct competencies rather than a
single capability \cite{li2023apibank}.

AgentBench extends evaluation beyond isolated tool calls by evaluating LLMs as
agents in multiple interactive environments. The study reports that poor long-term
reasoning, decision-making, and instruction following remain important limitations \cite{liu2023agentbench}.

These results are directly relevant to chained execution. A chain requires the model to
maintain a coherent objective across multiple decisions rather than merely produce
one valid API call.

\subsection{Security of Tool-Using Agents}

Tool use also creates security vulnerabilities because tool outputs may contain
information that the model interprets as instructions. InjecAgent specifically studies
indirect prompt injection in tool-integrated LLM agents and reports that external
content can manipulate agents into harmful actions or private-data exfiltration \cite{zhan2024injecagent}.

AgentDojo provides a dynamic environment for evaluating prompt-injection attacks
and defenses against agents operating on untrusted data. It includes realistic tasks
involving email, banking, and travel-booking environments and demonstrates that
modern agents can fail even without adversarial attacks, while prompt injections can
further compromise security properties \cite{debenedetti2024agentdojo}.

Agent Security Bench (ASB) expands security evaluation across multiple attack and
defense categories and reports vulnerabilities involving system prompts, user
prompts, tool use, and memory retrieval. Its results reinforce the observation that
agent security cannot be treated as a problem confined to the underlying language
model \cite{zhang2025asb}.

\section{Main Failure Modes in Chained Tool Calling}
\label{sec:failure-modes}

\subsection{Planning and Decomposition Failure}

A chained agent must transform a high-level objective into a sequence of executable
operations. Planning can fail when the model chooses an unnecessary step, omits a
required dependency, misunderstands the order of operations, or commits
prematurely to an incorrect strategy.

For example, suppose an agent is asked to analyze financial records and prepare a
report. It might first retrieve transactions, then classify them, then calculate totals,
then generate a report. If the classification procedure is wrong, subsequent
calculations can be internally consistent while still being based on an incorrect
interpretation.

The danger is therefore structural consistency without semantic correctness.
Later steps may appear successful because they correctly process the outputs of
earlier steps even though those outputs were wrong.

\subsection{Tool-Selection Failure}

An agent may select the wrong tool even when an appropriate tool exists.
Tool-selection failure can result from ambiguous descriptions, overlapping
capabilities, poor retrieval of tool documentation, or an incorrect interpretation of
the user's objective.

API-Bank's explicit separation of planning, retrieving, and calling
demonstrates that selecting and using tools are distinct competencies.

In a chained setting, wrong tool selection can be particularly damaging because the
output may still look plausible. A search tool, database query, or estimation API may
return a valid response even though it was not the appropriate source for the
intended task.

\subsection{Parameter and Schema Errors}

Even when the correct tool is selected, the model must generate valid parameters.
Errors may involve incorrect field names, invalid values, wrong units, omitted
constraints, or inappropriate defaults.

Gorilla explicitly identifies accurate input-argument generation as a major challenge in API calling.

A parameter error can also become a propagation problem. Consider a chain:
\begin{equation}
T_1(x) \rightarrow y
\label{eq:t1}
\end{equation}
\begin{equation}
T_2(y) \rightarrow z
\label{eq:t2}
\end{equation}
If $x$ is incorrectly specified, $y$ may be wrong but syntactically valid. The
second tool may therefore accept $y$, producing $z$ without raising an error.
The chain has then converted a detectable local mistake into a less detectable
downstream error.

\subsection{Error Propagation}

Error propagation is the defining reliability problem of chained execution.
Suppose each independent tool call has probability $p$ of being correct. Under a
simplified independence assumption, the probability that all $n$ calls are correct is:
\begin{equation}
P(\text{all correct}) = p^n
\label{eq:allcorrect}
\end{equation}
Thus, if every step has a 95\% probability of being correct, ten independent
steps would have an all-correct probability of approximately:
\begin{equation}
0.95^{10} \approx 0.599
\label{eq:example}
\end{equation}
This calculation is only illustrative because real agent failures are neither
independent nor identically distributed. Nevertheless, it demonstrates why long
chains can be substantially less reliable than individual operations.

More importantly, errors can be dependent. An early incorrect assumption may
systematically influence several subsequent decisions, causing correlated failures
rather than independent ones. This means that simply multiplying per-step success
rates may underestimate real-world risk.

\subsection{Misinterpretation of Tool Outputs}

Tool outputs are not automatically equivalent to verified facts. A tool can
return incomplete, stale, ambiguous, malformed, or context-dependent
information. The agent must interpret the output before using it in subsequent
actions.

A particularly dangerous pattern occurs when the agent treats a tool response as
authoritative even though the tool only returned a partial result. The resulting
assumption can then become embedded in the next tool call.

This problem is related to the broader distinction between data and instructions. In
tool-integrated systems, external content may contain text that the model interprets
as an instruction rather than merely as data. AgentDojo and InjecAgent demonstrate
why this distinction is difficult for current agents.

\subsection{State-Tracking and Context Drift}

Long chains require the agent to maintain an accurate representation of:
\begin{itemize}[noitemsep]
    \item the original user objective;
    \item completed actions;
    \item current external state;
    \item assumptions made during execution;
    \item information obtained from tools;
    \item constraints and permissions; and
    \item outstanding subgoals.
\end{itemize}

As the chain becomes longer, the possibility of state inconsistency increases. An
agent may forget an earlier constraint, confuse an intermediate value with a final
value, or continue executing after an external state has changed.

Agent evaluation research identifies long-term reasoning and decision-making as
important obstacles for usable agents, making state management a central
concern for autonomous chains.

\subsection{Prompt Injection and Instruction Contamination}

One of the most serious failure modes occurs when untrusted tool output contains
instructions designed to manipulate the agent. For example:
\begin{enumerate}[noitemsep]
    \item The user asks the agent to summarize documents.
    \item The agent calls a document-retrieval tool.
    \item One retrieved document contains malicious instructions.
    \item The model interprets those instructions as part of its operational context.
    \item The agent calls another tool to disclose information or perform an unauthorized action.
\end{enumerate}

The important characteristic is that the attacker may not need direct access to the
agent's system prompt. The malicious content can enter through a legitimate tool.

InjecAgent found substantial vulnerability to indirect prompt injection in
tool-integrated agents, while AgentDojo provides a broader environment for
evaluating this class of attack.

\subsection{Privilege and Capability Escalation}

Chained agents can combine individually harmless capabilities into a powerful sequence.
For example, one tool may read files, another may send messages, and a third may
access an external service. None of these capabilities necessarily creates
catastrophic risk by itself. However, chaining them can produce a path from
information retrieval to unauthorized disclosure or external action.

This suggests that security should evaluate not only individual tool permissions but also compositions
of permissions.

ASB's evaluation of attacks across system prompts, tool usage, and memory
retrieval illustrates why agent security must consider the entire workflow rather than
only isolated model responses.

\subsection{Irreversible and High-Impact Actions}

A major distinction between informational and operational tools is reversibility.

Examples of relatively reversible actions include:
\begin{itemize}[noitemsep]
    \item retrieving public information;
    \item calculating a value;
    \item generating a draft.
\end{itemize}

Examples of potentially irreversible or high-impact actions include:
\begin{itemize}[noitemsep]
    \item deleting files;
    \item transferring funds;
    \item sending messages;
    \item changing account permissions;
    \item modifying production infrastructure;
    \item publishing information.
\end{itemize}

An autonomous chain that proceeds directly from reasoning to irreversible
execution has a substantially different risk profile from one that pauses before
high-impact operations.

\subsection{Recovery Failure}

Agents do not only need to perform actions correctly; they need to recover when actions fail.
A weak recovery policy may:
\begin{enumerate}[noitemsep]
    \item repeat the same failed call;
    \item modify parameters without understanding the cause;
    \item switch tools unnecessarily;
    \item conceal the failure in the final answer; or
    \item continue executing despite an uncertain state.
\end{enumerate}

A robust agent should instead distinguish between recoverable and non-recoverable failures and know
when to stop.

This is especially important because automatic retry mechanisms can amplify
rather than reduce risk. Repeating a failed read operation may be harmless,
whereas repeatedly attempting a state-changing operation may produce multiple
unwanted effects.

\section{Why the Absence of Human Checkpoints Matters}

Human checkpoints can be understood as control barriers in the agent's execution
pipeline. Without a checkpoint:
\begin{equation}
\text{Plan} \rightarrow \text{Tool}_1 \rightarrow \text{Tool}_2
\rightarrow \text{Tool}_3 \rightarrow \text{Action}
\label{eq:nocheckpoint}
\end{equation}

With checkpoints:
\begin{equation}
\text{Plan} \rightarrow \text{Tool}_1 \rightarrow
\boxed{\text{Review}} \rightarrow \text{Tool}_2 \rightarrow
\boxed{\text{Review}} \rightarrow \text{Action}
\label{eq:checkpoint}
\end{equation}

The checkpoint provides an opportunity to detect incorrect assumptions
before they propagate.

However, human oversight is not necessarily required after
every operation. Excessive intervention can reduce efficiency
and undermine the purpose of autonomous agents. A more
useful model is risk-adaptive oversight, in which
checkpoints are triggered by properties such as:
\begin{itemize}[noitemsep]
    \item high uncertainty;
    \item sensitive data;
    \item unusual tool combinations;
    \item irreversible actions;
    \item large financial or operational consequences;
    \item privilege changes;
    \item conflicts between tool outputs;
    \item anomalous behavior;
    \item deviations from the original plan.
\end{itemize}

This approach treats human oversight as a scarce safety
resource rather than as a universal requirement.

\section{Current Approaches and Developments}

\subsection{Structured Tool Interfaces}

Structured APIs and explicit schemas can reduce
parameter-generation errors by constraining the model's
output space. API-Bank and related tool-learning research
demonstrate the importance of evaluating tool selection and
calling as separate capabilities.

\subsection{Retrieval-Augmented Tool Documentation}

Gorilla demonstrates that retrieving current API
documentation can improve the reliability of API calls and
reduce problems caused by changing tool interfaces.

\subsection{Sandboxing and Simulation}

ToolEmu proposes using an LM-emulated sandbox to
identify risks associated with tool-using agents without
requiring every real-world tool to be implemented. This
provides a scalable approach to testing potentially
dangerous actions before deployment \cite{ruan2024toolemu}.

\subsection{Adversarial Evaluation}

AgentDojo and InjecAgent demonstrate the importance of
testing agents against indirect prompt injection and untrusted
external content rather than evaluating only benign tasks.
ASB similarly evaluates a broad range of attacks and defenses
across different agent scenarios, emphasizing that agent
security involves multiple layers.

\subsection{Runtime Monitoring}

Runtime monitoring can inspect tool calls before execution,
compare actions with policy, detect anomalous sequences,
and prevent unauthorized operations. Such mechanisms are
especially important for chains because the safety of an
individual action may depend on what happened earlier.

Recent research increasingly emphasizes runtime controls and
capability restrictions rather than relying solely on
natural-language instructions to the model.

\section{Research Challenges and Limitations}

A major limitation of current research is that many
evaluations focus on task completion rather than the detailed
reliability of the execution process. A system can achieve a
successful final answer while making unsafe intermediate
decisions, particularly when the benchmark does not expose
those decisions as separate evaluation targets.

Recent evaluation surveys identify a gap between benchmark
performance and real deployment requirements, including
safety, robustness, cost, maintainability, and workflow
integration \cite{yehudai2025survey}.

Another challenge is the absence of standardized metrics
for cumulative chain reliability. Existing metrics often
report task success, tool-call accuracy, attack success, or
similar measures separately. What is needed is a
framework that records where in a chain an error occurred,
whether it propagated, whether recovery was attempted,
and whether a checkpoint could have prevented the final
failure.

Security evaluation also remains difficult because attacks
and defenses evolve rapidly. Agent Security Bench and
AgentDojo demonstrate that adversarial evaluation must
account for multiple attack surfaces and adaptive behavior.

Finally, human oversight itself requires careful design. A poorly
designed checkpoint may merely ask a human to approve an
action without providing enough context to understand the
decision. Effective oversight therefore requires concise
explanations, relevant evidence, action previews, and clear
indications of uncertainty.

\section{Research Gaps and Future Directions}

\subsection{Chain-Level Reliability Metrics}

Future benchmarks should report not only whether an agent completed a task but also:
\begin{itemize}[noitemsep]
    \item number of tool calls;
    \item number of retries;
    \item first point of failure;
    \item error propagation length;
    \item recovery success;
    \item irreversible actions attempted;
    \item policy violations;
    \item human interventions required.
\end{itemize}

This would allow researchers to distinguish a short successful trajectory from a
fragile trajectory that happens to reach the correct result.

\subsection{Provenance-Aware State Management}

Agents should maintain explicit provenance for
important intermediate values. Instead of
storing only: ``Customer balance = $X$'', the system should retain:
``Customer balance = $X$, obtained from Tool A at time $T$, using
query $Q$, with confidence/validation status $V$.'' Such provenance would
make it easier to identify stale or unverified assumptions before they enter
subsequent tool calls.

\subsection{Adaptive Checkpoint Placement}

Rather than fixed checkpoints, future systems could estimate the risk of continuing
autonomously. A checkpoint could be triggered when the predicted expected loss of
autonomous continuation exceeds the cost of human review.

Conceptually:
\begin{equation}
\text{Request human review if } \quad
P(\text{harm} \mid s,a) \times C_{\text{harm}} > C_{\text{review}}
\label{eq:checkpointrule}
\end{equation}
where $s$ is the current state, $a$ is the proposed action,
$C_{\text{harm}}$ is the estimated consequence of failure, and
$C_{\text{review}}$ is the cost of human intervention.

This creates a principled basis for balancing autonomy and oversight.

\subsection{Compositional Security Analysis}

Security testing should examine combinations of tools rather than only individual tool
permissions. A tool chain should be evaluated as a graph in which nodes represent
capabilities and edges represent possible transitions between capabilities.

This could reveal dangerous compositions that are
invisible when every tool is evaluated independently.

\subsection{Failure-Aware Agent Architectures}

Future agents should explicitly represent uncertainty and execution state. A robust
agent should be able to classify an intermediate result as:
\begin{itemize}[noitemsep]
    \item verified;
    \item plausible but unverified;
    \item contradictory;
    \item stale;
    \item malformed;
    \item unsafe to act upon.
\end{itemize}

This would help prevent the common failure pattern in which an uncertain
observation silently becomes a definitive premise for the next action.

\section{Conclusion}

Chained tool-calling agents represent an important transition from language
generation toward autonomous interaction with external systems. ReAct and
Toolformer established important foundations for reasoning and tool use, while later
benchmarks such as API-Bank, AgentBench, AgentDojo, InjecAgent, ToolEmu, and
Agent Security Bench have exposed limitations in planning, API invocation,
long-horizon reasoning, security, and robustness.

The central problem with autonomous chains is not simply that an agent can make a
mistake. It is that an early mistake can become the basis for later decisions.
Incorrect plans, tool choices, parameters, interpretations, or state assumptions can
therefore propagate across multiple steps. When the chain includes high-impact or
irreversible operations, the resulting failure can move from an informational error to a
real-world consequence.

Human checkpoints provide an effective mechanism for interrupting this
propagation, but continuous human approval is neither scalable nor desirable.
The more promising direction is adaptive oversight combined with structured tool
interfaces, capability restrictions, provenance tracking, runtime monitoring,
sandboxing, adversarial testing, and explicit uncertainty management.

Ultimately, safe autonomous tool use requires systems to know not only how to call
a tool, but also when not to call it, when to verify its output, when to recover,
and when to stop. Future research should therefore evaluate agents as complete
socio-technical execution systems rather than as language models that merely
produce successful final answers.

\bibliographystyle{plain}
\bibliography{references}

\end{document}
